% Optional corollary to gaussian-simple-targets.tex.
\begin{corollary}
For every fixed $d,N,L$, the globally margin-conditioned Gaussian function
prior converges, as $\sigma\downarrow0$, to the uniform distribution on
$\mathcal M_d=\{\pm1,\pm x_1,\ldots,\pm x_d\}$.
For data labeled by any $f\in\mathcal M_d$, the posterior converges to the
uniform distribution on the members of $\mathcal M_d$ consistent with the
data. In fact this posterior convergence is uniform over all such labeled
samples, including arbitrary sample sizes.
For uniform IID inputs, if
\[
 n\ge\log_2\frac{2d+1}{\delta},
\]
then with probability at least $1-\delta$ the limiting posterior is the point
mass at $f$.
\end{corollary}

\begin{proof}
The norm theorem shows that the prior mass outside $\mathcal M_d$ tends to
zero relative to the mass of any one member. Their masses are exactly equal
at every positive variance. Indeed, swapping columns $0$ and $i$ in the first
weight matrix preserves its independent Gaussian law. Let $T_i$ be the
involution of the Boolean cube defined by
\[
 (T_i x)_i=x_i,\qquad (T_i x)_j=x_i x_j\quad(j\ne i).
\]
Since the entire network, as a function of its augmented input, is odd, this
column swap transforms its output into
\[
 z_{w'}(x)=x_i z_w(T_i x).
\]
It therefore bijects the strict-margin implementation sets for $+1$ and
$+x_i$. Negating the readout preserves its Gaussian law and supplies the
negative functions. Thus all $2(d+1)$ implementation sets have equal mass.

Conditioning on any sample labeled by $f\in\mathcal M_d$ retains at least
$f$. If $\pi$ is the function prior, the posterior mass outside
$\mathcal M_d$ is bounded by $\pi(\mathcal M_d^c)/\pi(f)$, independently of
the sample and its size. Within $\mathcal M_d$, consistency conditioning is
exactly uniform. This proves the posterior convergence claim.
Under uniform inputs, each distinct $g\in\mathcal M_d$ disagrees with $f$ on
at least half the cube. Hence its probability of surviving $n$ labels is at
most $2^{-n}$. A union bound over the at most $2d+1$ competitors proves the
last assertion.
\end{proof}
